\begin{definition}[Regular half-edge coordinates and closed coefficients]%
    \label{def:peps_regular_half_edge_closed}
    \lean{TNLean.PEPS.RegularHalfEdgeConfig,
        TNLean.PEPS.regularHalfEdgeLeftMul,
        TNLean.PEPS.regularHalfEdgeRightMul,
        TNLean.PEPS.regularHalfEdgeOperators,
        TNLean.PEPS.regularHalfEdgeOfBonds,
        TNLean.PEPS.regularHalfEdgeOfOperators,
        TNLean.PEPS.regularHalfEdgeConfigEquivOperatorsBonds,
        TNLean.PEPS.fullRegionHalfEdgeConfigEquiv,
        TNLean.PEPS.regularProjectorClosedState,
        TNLean.PEPS.regularHalfEdgeOfBonds_tail,
        TNLean.PEPS.regularHalfEdgeOfBonds_head,
        TNLean.PEPS.regularHalfEdgeOperators_ofBonds,
        TNLean.PEPS.regularHalfEdgeOperators_ofOperators,
        TNLean.PEPS.regularHalfEdgeOperators_rightMul,
        TNLean.PEPS.regularHalfEdgeOperators_leftMul,
        TNLean.PEPS.regularHalfEdgeOfBonds_eq_rightMul}
    \leanok
    \uses{def:peps_regular_projector_twisted_region}
    Let $G$ be a finite group and let $\Gamma=(V,E)$ be a finite
    simple graph, with each edge oriented from $s(e)$ to $t(e)$.
    A half-edge configuration assigns $\alpha_{v,e}\in G$ to
    every incident vertex-edge pair. Vertex and bond translations are
    $(k\alpha)_{v,e}=k_v\alpha_{v,e}$ and
    $(\alpha r)_{v,e}=\alpha_{v,e}r_e$, respectively.
    Define the operator quotient
    $u(\alpha)_e=\alpha_{t(e),e}\alpha_{s(e),e}^{-1}$.
    For $u,\eta\in G^E$, let $\alpha(u,\eta)$ have tail labels
    $\eta_e$ and head labels $u_e\eta_e$. The map
    $(u,\eta)\mapsto\alpha(u,\eta)$ is a bijection, with inverse
    $\alpha\mapsto(u(\alpha),(\alpha_{s(e),e})_e)$.
    Write $\alpha(u)=\alpha(u,1)$. Then
    $\alpha(u,\eta)=\alpha(u)\eta$,
    $u(\alpha r)=u(\alpha)$, and
    $u(k\alpha)_e=k_{t(e)}u(\alpha)_ek_{s(e)}^{-1}$.

    With $P_v$ the regular averaging projector, the canonical closed
    coefficient vector with inserted operators $u$ is
    \begin{align}
        C_u(\alpha)
        =\sum_{\eta\in G^E}\prod_{v\in V}
        P_v\bigl(\alpha_v,\alpha(u,\eta)_v\bigr).
        \label{eq:peps_half_edge_closed}
    \end{align}
    These are the full-region contractions of
    Definition~\ref{def:peps_regular_projector_twisted_region}.
\end{definition}

\begin{theorem}[Supported expansion of invariant half-edge functions]%
    \label{thm:peps_regular_invariant_state_span}
    \lean{TNLean.PEPS.eq_sum_regularProjectorClosedState_of_invariant,
        TNLean.PEPS.eq_sum_supported_regularProjectorClosedState_of_invariant,
        TNLean.PEPS.mem_span_supported_regularProjectorClosedState_of_invariant,
        TNLean.PEPS.regularProjectorClosedState_eq_zero_of_not_support}
    \leanok
    \uses{def:peps_regular_half_edge_closed}
    Suppose $f(k\alpha)=f(\alpha r)=f(\alpha)$ for every vertex
    translation $k$, bond translation $r$, and half-edge configuration
    $\alpha$. Let $\mathcal P\subseteq G^E$ satisfy
    $u(\alpha)\notin\mathcal P\Rightarrow f(\alpha)=0$. Then
    \begin{align}
        f=\sum_{u\in\mathcal P} f\bigl(\alpha(u)\bigr)C_u.
        \label{eq:peps_supported_orbit_expansion}
    \end{align}
    If $\mathcal P$ is preserved by the endpoint gauges
    $u_e\mapsto k_{t(e)}^{-1}u_ek_{s(e)}$, then every $C_u$ with
    $u\in\mathcal P$ also vanishes when $u(\alpha)\notin\mathcal P$.
    This is the finite regular-coordinate calculation associated with
    the closure argument of \cite[Theorem~5.5]{Schuch2010PEPS}.
\end{theorem}

\begin{proof}\leanok
    \uses{def:peps_regular_half_edge_closed}
    Put $\kappa=|G|$. The independent vertex projectors have kernel
    \begin{align}
        \prod_vP_v(\alpha_v,\beta_v)
        =\kappa^{-|V|}\sum_{k\in G^V}
        \mathbf1_{\{\alpha=k\beta\}}.
        \label{eq:peps_half_edge_average}
    \end{align}
    Summing $f(\beta)$ against (\ref{eq:peps_half_edge_average}) gives
    $\kappa^{-|V|}\sum_k f(k^{-1}\alpha)=f(\alpha)$.
    Reindex $\beta$ by the bijection $(u,\eta)\mapsto\alpha(u,\eta)$.
    Since $\alpha(u,\eta)=\alpha(u)\eta$, right invariance replaces
    $f(\alpha(u,\eta))$ by $f(\alpha(u))$. Summing over $\eta$ gives
    (\ref{eq:peps_half_edge_closed}). If $u\notin\mathcal P$,
    the coefficient $f(\alpha(u))$ is zero because
    $u(\alpha(u))=u$. Removing these zero coefficients gives
    (\ref{eq:peps_supported_orbit_expansion}).

    A nonzero indicator in (\ref{eq:peps_half_edge_average}), with
    $\beta=\alpha(u,\eta)$, implies
    $u(\alpha)_e=k_{t(e)}u_ek_{s(e)}^{-1}$.
    Gauge invariance of $\mathcal P$ therefore forces every indicator
    to vanish outside its prescribed support.
\end{proof}

\begin{theorem}[Physical contraction of the regular closed coefficients]%
    \label{thm:peps_regular_closed_physical_contraction}
    \lean{TNLean.PEPS.regularProjectorClosedState_eq_twistedRegionMatrix,
        TNLean.PEPS.sum_regularProjectorClosedState_eq_stateCoeff}
    \leanok
    \uses{def:peps_regular_half_edge_closed,
        thm:peps_regular_twisted_region_accessibility}
    Let the site maps $A_v$ be regular $G$-injective. For each physical
    configuration $\sigma$, the inserted physical PEPS coefficient is
    \begin{align}
        \sum_\alpha\left(\prod_v A_{v,\sigma_v}(\alpha_v)\right)C_u(\alpha)
        =\sum_{\eta\in G^E}\prod_v
        A_{v,\sigma_v}\bigl(\alpha(u,\eta)_v\bigr).
    \end{align}
\end{theorem}

\begin{proof}\leanok
    \uses{thm:peps_regular_open_region_accessibility}
    Expand (\ref{eq:peps_half_edge_closed}), exchange the two finite sums,
    and factor the sum over independent half-edge labels into one sum
    per vertex. At each vertex use $A_vP_v=A_v$. The remaining sum
    over the shared labels $\eta$ is the actual inserted contraction.
\end{proof}
